\section{Joint Residual Refinement}\label{app:joint-refinement-v92}
The residual-energy first stage identifies small certifying spans. When its sufficient tail exceeds the requested precision, we refine the original completion domain with a reference-residual certificate. The two stages share the original candidate, ridge penalty, completion boxes and full-feature retrained comparator.

\begin{proposition}[Joint reference-residual bound]\label{prop:joint-refinement-v92}
Let $G_0\succ0$ be a uniform lower bound on $M_X$, and verify $T^\top G_0T\succeq\mu I$ for an invertible actual matrix $T$ and $\mu>0$. For any fixed reference $\vartheta$, put $g_\vartheta(X)=M_X\vartheta-r_X$. Then
\begin{align}
h_\beta(X)&=J_X(\beta)-J_X(\vartheta)\nonumber\\
 &\quad+g_\vartheta(X)^\top M_X^{-1}g_\vartheta(X),\label{eq:reference-identity-v92}\\
h_\beta(X)&\le P_\vartheta(X):=J_X(\beta)-J_X(\vartheta)\nonumber\\
 &\quad+\mu^{-1}\|T^\top g_\vartheta(X)\|^2.\label{eq:joint-refinement-v92}
\end{align}
On an original completion subbox, $P_\vartheta$ is a polynomial of degree at most four in normalized cell coordinates $z\in[0,1]^d$.
\end{proposition}
\begin{proof}
Completing the square around the exact full-feature ridge solution gives~\eqref{eq:reference-identity-v92}. The verified metric implies $M_X^{-1}\preceq\mu^{-1}TT^\top$, which gives~\eqref{eq:joint-refinement-v92}. Each row is affine in its uncertain cells. The objective difference is quadratic, $g_\vartheta$ is quadratic, and its squared transformed norm has degree at most four.
\end{proof}

The polynomial retains the dependence between the objective difference and the normal-equation residual. Write $x_{ij}=x^0_{ij}+d_a z_a$ for cell $a=(i,j)$, with $x^0$ an endpoint table and signed length $d_a$. Before applying $T^\top$, the constant gradient coefficient is $g_\vartheta(X^0)$; the linear coefficient is $w_i d_a[(x_i^{0\top}\vartheta-y_i)e_j+\vartheta_j x_i^0]$. The square coefficient is $w_i d_a^2\vartheta_j e_j$. For two distinct cells $a=(i,j)$ and $b=(i,k)$ in one row, their product coefficient is $w_i d_a d_b(\vartheta_k e_j+\vartheta_j e_k)$. Squaring these explicit coefficients constructs the joint polynomial with outward arithmetic.

\paragraph{Lifted row residuals and box localizers.}
For a row with several uncertain entries, introduce a normalized residual $u_i\in[0,1]$ through
\[
 x_i^\top\vartheta-y_i=\ell_i+s_i u_i,
\]
where $s_i>0$ and $[\ell_i,\ell_i+s_i]$ contains the exact residual range. We use dyadic endpoints and spans. Two outward-rounded affine inequalities encode this relation, together with their products by each normalized variable and its complement. The transformed normal residual is bilinear in cell coordinates and the corresponding $u_i$. Its moment basis therefore has at most $1+2d+m$ entries for $d$ uncertain cells in $m$ rows.

Let $v_i$ contain the constant, the normalized cells in row $i$, and $u_i$. For each coordinate $q$ in that row, we add the localizing constraint
\begin{equation}
 \mathcal L_{i,q}(y)
 :=\mathbb E_y\!\left[q(1-q)v_i v_i^\top\right]\succeq0.
 \label{eq:row-localizer-v92}
\end{equation}
Here $\mathbb E_y$ denotes substitution of monomials by lifted moments. Every actual completion satisfies~\eqref{eq:row-localizer-v92}, since $q(1-q)\ge0$ and $v_i v_i^\top\succeq0$. For a row containing $k_i$ uncertain cells, this adds $k_i+1$ matrices of order $k_i+2$. These constraints retain the row's box geometry in the fourth-degree relaxation.

\paragraph{A validated moment dual.}
We use polynomial moment relaxations~\citep{lasserre2001}, modeled in CVXPY~\citep{diamond2016cvxpy}, with SCS numerical proposals~\citep{odonoghue2016scs}.
Let $y_\alpha=z^\alpha$ for every required monomial and its factors, with $0\le y_\alpha\le1$. Product relations give $y_{\alpha+\gamma}\le y_\alpha,y_\gamma$ and $y_{\alpha+\gamma}\ge y_\alpha+y_\gamma-1$. Square tangents use $y_{2\alpha}\ge2t y_\alpha-t^2$ for $t\in\{0,1/4,1/2,3/4,1\}$. Write these inequalities as $Ay\le b$. The moment matrix $\mathcal M(y)$ indexed by the constant, linear and within-row quadratic gradient monomials satisfies $\mathcal M(y)\succeq0$ on every genuine completion.

Enclose the polynomial coefficients above as $c_0,c$. For any nonnegative multiplier $\lambda$ and any matrix $R$, write
\[
 \langle RR^\top,\mathcal M(y)\rangle=s_0+s^\top y.
\]
Every allowed completion obeys the independently checkable upper bound
\begin{equation}
 P_\vartheta(X)\le c_0+\lambda^\top b+s_0
       +\sum_j\max\{0,c_j-(A^\top\lambda)_j+s_j\}.
 \label{eq:validated-moment-dual-v92}
\end{equation}
Indeed, add the nonnegative moment quadratic and the nonnegative inequality slack to the polynomial, then bound the remaining linear form on $[0,1]$. Numerical conic optimization proposes the multipliers. The implementation rounds $\lambda$ to nonnegative dyadics and stores a dyadic factor $R$, which establishes positive semidefiniteness directly. Outward evaluation of~\eqref{eq:validated-moment-dual-v92} includes the complete dual residual. The independent auditor reconstructs the original polynomial, actual metric, multiplier constraints and factor products as exact rationals.

For the localized relaxation, each matrix in~\eqref{eq:row-localizer-v92} receives its own dyadic factor. Their inner products contribute additional constant and monomial coefficients to $s_0$ and $s$ in~\eqref{eq:validated-moment-dual-v92}. The same residual correction proves the resulting upper bound. The independent checker also reconstructs the residual range, affine links, and every localizing polynomial $q-q^2$.

\paragraph{Reference and subdivision rules.}
A fixed 64-start feasible search supplies a completed-table lower witness. An eight-step finite-pool minimax exchange proposes a reference, using 32 starts per exchange step and fixed random seeds. Each proposed reference receives the uniform polynomial certificate above. For an unresolved node, we bisect the cell maximizing relaxed moment variance times its polynomial coefficient sensitivity. The parent remains active until both covering children finish. The largest active upper, intersected with the first-stage upper, and the best feasible original-table lower give the global interval. Every reference fit, witness search, polynomial construction, conic solve and dual evaluation enters the refinement clock.

The prespecified new-mask computation cap is 600 seconds, with completed-evidence checkpoints at 10, 60, 120, 300 and 600 seconds. All twenty new masks finish within 76.21 seconds. The original rank-three guarantee belongs to the structural first stage; this completion-domain refinement supplies additional accuracy when needed. Its computational cost depends on the unresolved cells and the moment relaxation.

\subsection{Row-edge duals and certificate screening}\label{app:edge-refinement-v22}
For a small number of dense uncertain rows, we use Theorem~\ref{thm:row} to place the coupling certificate on a product of row edges. Let $m$ be the number of uncertain rows, fix a reference $\vartheta$, and write
\begin{align*}
 g_0&=G_0\vartheta-r_0,\\
 \gamma_i(x_i)&=T^\top\!\left[w_i x_i(x_i^\top\vartheta-y_i)+g_0/m\right].
\end{align*}
Then $\sum_i\gamma_i=T^\top g_\vartheta(X)$. With
\[
 f_i(x_i)=w_i[(x_i^\top\beta-y_i)^2-(x_i^\top\vartheta-y_i)^2],
\]
the joint majorant has the form
\[
 P_\vartheta(X)=c+\sum_i f_i(x_i)+\mu^{-1}\Big\|\sum_i\gamma_i(x_i)\Big\|^2,
\]
where $c=J_0(\beta)-J_0(\vartheta)$ contains the complete rows and ridge term.

\begin{proposition}[Row-edge coupling certificate]\label{prop:edge-dual-v22}
Let $Q\succeq0$ have scalar block $Q_{00}=\alpha$, blocks $Q_{i0}=d_i$, diagonal blocks $Q_{ii}=H_i-\mu^{-1}I$, and off-diagonal row blocks $Q_{ij}=-\mu^{-1}I$. Define
\[
 R_i=f_i+\gamma_i^\top H_i\gamma_i+2d_i^\top\gamma_i,
 \qquad b_i\ge\max_{x_i\in E_i}R_i(x_i),
\]
where $E_i$ is the union of the row box's edges. Then
\[
 H^*\le U:=c+\alpha+\sum_i b_i.
\]
\end{proposition}
\begin{proof}
For $v=(1,\gamma_1,\ldots,\gamma_m)$, direct expansion gives
\begin{align*}
 U-P_\vartheta(X)&=v^\top Qv+\sum_i[b_i-R_i(x_i)]\\
 &\ge0\quad\text{on }\prod_iE_i.
\end{align*}
The row-face theorem places an original-regret maximizer in this product, and the joint majorant bounds its value.
\end{proof}

Each row feature is quadratic in its normalized cells. On one edge, $R_i$ is a quartic in one scalar. The numerical exchange method solves the coupling SDP on a finite set of row points and adds distinct violating points from the row oracles. Its reported upper uses the complete row-oracle bounds. The verifier reconstructs the exact polynomial coefficients and cross-row matrix blocks, checks a rational congruence of $Q$ with positive diagonal-dominance margins, and encloses every quartic maximum by integer Bernstein subdivision.

The certificate also restricts the remaining search. For a target value $\tau<U$, every edge-domain completion with original regret exceeding $\tau$ must satisfy
\begin{equation}
 R_i(x_i)>b_i-(U-\tau)\quad\text{for every uncertain row }i.
 \label{eq:edge-screen-v22}
\end{equation}
Indeed, a violation of one inequality makes the corresponding nonnegative slack at least $U-\tau$. An exact Bernstein upper below this cutoff excludes the entire edge segment. The retained segments cover every possible target violation.

We partition this cover using the fixed-reference residuals $x_i^\top\vartheta-y_i$. A residual slab clips each edge to an interval, so the same quartic verifier applies. The numerical master proposes the row with largest residual variance and a dyadic split at its weighted mean, clipped to the inner 70\% of the available range. Both complementary children are verified before replacing their parent. The final original-regret upper is
\[
 \max\{\tau,\max_{B\in\mathcal L}U_B\},
\]
intersected with the verified root bound; $\mathcal L$ is the complete active cover. Every node retains the original candidate, feature intervals, ridge penalty and full-feature comparator.

The dense-row protocol fixes the target $\varepsilon=10^{-8}$, uses $\tau=L+0.99\varepsilon$, and sets the complete cap at 1800 seconds. It uses a 4096-start feasible search, an eight-step reference exchange with 32 starts per step, up to 180 seconds for the root proposal and 90 seconds for each subsequent proposal, and at most 65 evaluated nodes. Exact bound evaluation, screening, reference construction and all completed child work enter the algorithm clock. A separate rational replay checks the original witness, matrix, exclusions and complete tree. The implementation and all four dense-row conditions use this common protocol.

\subsection{Algebraic Row Compression and Observed-Column Elimination}\label{app:compact-profile}
The general-coarsening solver combines two exact reductions with the validated row-edge bound. Algebraic row coordinates reduce the coupling matrices, and analytic elimination of fully observed columns reduces the outer comparator search when uncertainty is concentrated in a few columns.

\paragraph{Algebraic row coordinates.}
For row $i$, let $x_i=x_i^0+D_i z_i$, where $D_i$ has one scaled coordinate vector for each of its $k_i$ uncertain cells and $z_i\in[0,1]^{k_i}$. Fix a reference $\vartheta$, put $e_i^0=x_i^{0\top}\vartheta-y_i$ and $q_i=\vartheta^\top D_i z_i$, and define
\[
 u_i(z_i)=\begin{bmatrix}q_i\\z_i(e_i^0+q_i)\end{bmatrix},\qquad
 V_i=w_i\begin{bmatrix}T^\top x_i^0&T^\top D_i\end{bmatrix}.
\]
The transformed row-gradient increment is exactly $V_i u_i$. Thus, for a constant $a_0$ containing the endpoint gradients, observed rows and ridge,
\[
 T^\top g_\vartheta(X)=a_0+\sum_i V_i u_i(z_i).
\]
This uses at most $k_i+1$ scalar coordinates per row, independently of the number of model features. Exact zero columns can be omitted, and diagonal dyadic rescaling preserves the identity.

Write the joint majorant as
\[
 C+\sum_i\bar f_i(z_i)+c\Big\|\sum_iV_i u_i(z_i)\Big\|^2,
 \qquad c=\mu^{-1},
\]
where $C$ absorbs $c\|a_0\|^2$ and $\bar f_i$ absorbs $2c\,a_0^\top V_i u_i$. Choose $H_i\succ0$, $d_i$ and $\alpha$ satisfying
\begin{equation}
 \begin{bmatrix}\alpha&0\\0&c^{-1}I\end{bmatrix}
 -\sum_i\begin{bmatrix}d_i^\top\\V_i\end{bmatrix}
 H_i^{-1}\begin{bmatrix}d_i&V_i^\top\end{bmatrix}\succeq0.
 \label{eq:compact-schur}
\end{equation}
Two Schur complements give
\[
 c\Big\|\sum_i V_i u_i\Big\|^2
 \le\alpha+\sum_i\{u_i^\top H_i u_i+2d_i^\top u_i\}.
\]
Consequently, the sum of the corresponding row-edge quartic maxima plus $C+\alpha$ bounds the original regret. The numerical master uses one coupling cone and one small block cone per row. The independent verifier checks each local inverse exactly, encloses its Schur contribution by a dyadic positive-semidefinite upper matrix, verifies the central matrix by rational LDL elimination, and checks all row quartics by Bernstein subdivision. These checks retain all cross-row terms.

\paragraph{Curvature from uncertain rows.}
Let $c_i,h_i$ be row midpoints and halfwidths, and let
\[
 D=\operatorname{diag}\!\left(\sum_{i\notin\mathcal A}w_i k_i h_{ij}^2\right)_j.
\]
For $0<t<1$, matrix Young's inequality gives the uniform lower matrix
\begin{equation}
 M_X\succeq G_t:=G_0+(1-t)\sum_{i\notin\mathcal A}w_i c_i c_i^\top
 -(t^{-1}-1)D.\label{eq:young-curvature}
\end{equation}
Indeed, $(c+\delta)(c+\delta)^\top\succeq(1-t)cc^\top-(t^{-1}-1)\delta\delta^\top$ and $\delta\delta^\top\preceq k_i\operatorname{diag}(h_i^2)$. The implementation uses $t=1/2$ when its rounded metric is verified positive. This strengthens the common curvature used by the reference certificate.

\begin{proposition}[Profiling fully observed columns]\label{prop:observed-profile}
Let $J$ contain every uncertain feature column, let $K$ be its complement, and put $A=\lambda I+X_K^\top W X_K\succ0$. For $t\in\mathbb R^{|J|}$, define
\[
 F_J(t)=\max_{X\in\mathcal X}\{J_X(\beta)-\min_b J_X(\theta_J=t,\theta_K=b)\}.
\]
For any reference $\vartheta$ with $\vartheta_J=t$, its fixed-completion objective equals
\begin{equation}
 J_X(\beta)-J_X(\vartheta)+g_{\vartheta,K}(X)^\top A^{-1}g_{\vartheta,K}(X).
 \label{eq:profile-identity}
\end{equation}
Each fixed-completion function is concave quadratic in $t$, with negative quadratic coefficient
\[
 S_X=M_{JJ}-M_{JK}A^{-1}M_{KJ}\preceq M_{JJ}.
\]
If $U_0\ge H^*$, then
\[
 H^*=\max_{|t_j-\beta_j|\le\sqrt{U_0/\lambda}}F_J(t).
\]
\end{proposition}
\begin{proof}
Completing the square in $b$ proves~\eqref{eq:profile-identity}; a second completion of the square gives $S_X$. The full optimum satisfies $\lambda\|\theta_X^*-\beta\|^2\le h_\beta(X)\le U_0$, proving the containing box. Maximization over that box and the original completion domain recovers the full retraining optimum.
\end{proof}
For each corner query, a verified factor of $A$ bounds the last term in~\eqref{eq:profile-identity}. The transformed gradient of an uncertain row varies along the fixed direction obtained from its observed entries, giving a single algebraic row coordinate. A maximizing pair $(X,b)$ at fixed $t$ admits a row-edge completion by Theorem~\ref{thm:row}. Thus the compressed row certificate applies to each query. Theorem~\ref{thm:interp}, with
\[
 H_j=\lambda+\sum_iw_i\max\{L_{ij}^2,U_{ij}^2\},
\]
certifies every outer box. This computational reduction retains unrestricted retraining in $K$. Its search dimension is $|J|$; the residual-energy theorem separately controls certification dimension through residual structure.

\paragraph{Choosing the conic representation.}
The implementation compares the sizes of the two formulations using only $p$ and the uncertain-cell counts $k_i$. Put
\[
 \begin{aligned}
 q&=1+\sum_i k_i+\sum_i k_i(k_i+1)/2,\\
 C_{\rm moment}&=q(q+1)/2,
 \end{aligned}
\]
and
\[
 C_{\rm row}=\frac{(p+1)(p+2)}2+
 \sum_i\frac{(p+k_i+2)(p+k_i+3)}2.
\]
These count symmetric entries in the principal moment cone and the compressed block cones. The solver chooses the moment formulation when $C_{\rm moment}<C_{\rm row}$ and compressed rows otherwise. This geometry-based rule uses the same original objective and certificate tolerance in both branches.

\paragraph{Common general-coarsening protocol.}
Every condition uses rank-one/two/three structural selection and a one-second initial search. Both refinement branches start with 64 feasible starts. The moment branch uses eight reference-exchange rounds with 32 starts, one root-reference update, 60-second root and 40-second child proposals, and at most 33 evaluated nodes. A leading relaxed covariance direction and its mean propose complementary halfspaces; their nonnegative linear slacks and products with normalized variables enter the validated dual. The compressed branch uses sixteen reference exchanges with 32 starts and a 90-second root proposal. Remaining compressed cases with at most three uncertain columns use the profiled comparator search with at most 1023 evaluated boxes; other cases use fixed-reference residual slabs with at most 513 nodes and 60 seconds per proposal. The complete cap is 9000 seconds. Representation selection, input construction, reference fitting, numerical proposals, exact evaluations and completed partial-child work enter the clock. A parent stays active until both covering children are verified. A separate rational replay checks the representation rule, original witness, every query and the full cover.

\begin{algorithm}[t]
\caption{Original-table regret certification}\label{alg:certificate}
\begin{algorithmic}[1]
\REQUIRE Row boxes, $y,w,\lambda,\beta$, rank cap, tolerance $\varepsilon$, total budget $B$, structural-stage budget $B_s\le B$
\STATE Start the complete clock; charge input and preprocessing to $B$.
\STATE Form $G_0,r_0$ from complete rows and ridge once.
\STATE Propose dyadic $T$; verify its metric and an original upper $U_0$.
\STATE Propose uncertainty-gradient spans or signal-row spans plus an anchor-correction direction; verify the actual Gram/projector.
\STATE Select a span using verified spectral, moment or structural energy budgets.
\STATE Establish a finite root; initialize a feasible original lower $L$.
\WHILE{complete time is below $B_s$ and $U-L>\varepsilon$}
\STATE Pop the leaf with largest upper; bisect its weighted widest side.
\STATE Query row maxima; bound both children; update a feasible table.
\STATE Retrain that table in all $p$ features and update $L$.
\STATE Update $U$ from surviving leaf uppers, tail and $U_0$.
\ENDWHILE
\STATE If $U-L>\varepsilon$, apply the validated joint residual refinement; use observed-column elimination or row-edge subdivision for the general-coarsening protocol (Appendix~\ref{app:compact-profile}).
\STATE \textbf{return} $[L,U]$, completed-table witness, error ledger and stop reason
\end{algorithmic}
\end{algorithm}

